\section{Related Work}
\label{sec:related}

\para{Adaptive attacks.} Automated red teaming began with attacker language models and large-scale
human red teaming~\cite{perez2022redteaming,ganguli2022redteaming} and with gradient-based universal
suffixes~\cite{zou2023gcg}, and is now standardised by HarmBench~\cite{mazeika2024harmbench}. Nearly
every method since iterates on failure.
PAIR~\cite{chao2023pair} refines a candidate using the target's refusal, TAP~\cite{mehrotra2023tap}
adds tree search with pruning, Rainbow Teaming~\cite{samvelyan2024rainbow} maintains a
quality-diversity archive over attack styles, GOAT~\cite{pavlova2024goat} picks a published attack
technique per turn, and curiosity-driven red teaming~\cite{hong2024curiosity} adds a novelty bonus
to an RL attacker. \citet{andriushchenko2024adaptive} state the thesis most directly: adaptive
attacks are necessary for accurate robustness evaluation, and no single method generalises across
targets. Their ablation on \llamaseven moves from 0\% to 50\% to 100\% attack success as random
search and then self-transfer initialisation are added, and their R2D2 result is the built-in
negative control: self-transfer gives only 12\% there, while switching to an in-context template
family gives 90\%. Best-of-$N$~\cite{hughes2024bon} is the canonical non-adaptive control, being
memoryless by construction, and establishes the power law $-\log \mathrm{ASR}(N) = aN^{-b}$ that
we use as the extrapolative bound in \secref{sec:results}. Unlike all of this work, we hold the
attack set fixed and vary the \emph{protocol} that schedules it under a fixed budget, so
adaptivity is not confounded with compute.

\para{Return curves and reproducibility.} Whether adaptive allocation can beat uniform allocation at
all depends on the shape of the return curve. \citet{halder2026scaling} report a polynomial regime
without prompt injection and an exponential regime with it, and \citet{topol2026process} apply
process mining to 8{,}575 scored red-teaming events to show that mutator effectiveness is asymmetric
across models and that time-to-jailbreak distributions differ by an order of magnitude. This is
independent 2026 confirmation of the lesson we build on, that which adaptation works is
model-specific. \citet{fang2026foundry} reproduce 30 published attacks in a unified harness to a
mean deviation of $+0.26$ percentage points, which is the kind of control that makes replay-based
comparisons like ours credible.

\para{Nulls, forecasting, and elicitation.} \citet{jones2025forecasting} give the closest prior
work: they replace a binary null with a continuous elicitation probability and extrapolate extreme
quantiles via a Gumbel-tail fit, forecasting risk two orders of magnitude beyond the observed
budget. Their final section uses forecasts to allocate red-teaming compute, which is the direct
ancestor of our portfolio experiment. They also propose adaptively stopping unpromising queries as
future work. We take that proposal seriously enough to test it, and find that the within-run
signal it would rely on is not diagnostic on the population that matters.

\para{Hidden robustness and false nulls.} A second literature supplies the phenomena.
\citet{wei2023jailbroken} give the taxonomy, splitting failures into competing objectives and
mismatched generalisation, which is in effect a principled axis set for a design space. Shallow
safety alignment~\cite{qi2024shallow} shows alignment mostly adapts the first few output tokens,
which explains why suffix, prefilling, and decoding attacks all work. Sleeper
agents~\cite{hubinger2024sleeper} persist through SFT, RL, and adversarial training, with
adversarial training teaching models to hide the trigger better. Fine-tuning breaks guardrails with
ten examples~\cite{qi2023finetuning}. \citet{lucki2024unlearning} recover unlearned capabilities
with attacks previously reported as ineffective, and \citet{che2025tampering} show that input-output
evaluations only lower-bound worst-case behaviour. On the other side, latent adversarial
training~\cite{sheshadri2024lat} and constitutional classifiers~\cite{sharma2025classifiers} are
rare cases where a null appears to be earned, the latter at a cost of over 3{,}000 human red-teaming
hours. \citet{betley2025emergent} add a measurement problem that compounds all of this: narrowly
fine-tuned models behave inconsistently and sometimes act aligned, so single-sample evaluation has
low power by construction. Whether adaptive search should pay depends on the dimensionality of the vulnerability
boundary, which is itself disputed: \citet{arditi2024refusal} report a single refusal direction,
while \citet{wollschlager2025geometry} report multi-dimensional concept cones. We treat this as a
moderator rather than settled, and our results are consistent with the family-dependent view.

\para{Statistics for evaluation.} \citet{miller2024errorbars} imports standard errors, clustering,
paired inference, and power analysis into evaluation, and shows clustered standard errors can be
$3.05\times$ larger than naive ones. \citet{bowyer2025clt} show CLT intervals badly understate
uncertainty below a few hundred datapoints, which matters here because adaptive protocols
concentrate budget on small subsets. Anytime-valid inference~\cite{ramdas2022savi,lindon2022anytime}
is the correct frame for data-dependent stopping, and adaptive item selection with preserved
coverage has been developed for capability estimation~\cite{polo2024tinybenchmarks,krumdick2026wild,wu2026faq,tolochinsky2026valid}.
\citet{dorner2024limits} bound how far an LLM judge can substitute for labels, which is why we
report both a judge-based and a rule-based outcome throughout. The red-teaming literature has the
optional-stopping problem and this literature has the fix; we join them, and find, surprisingly,
that the fix addresses the smaller of two errors.

\para{Bayesian experimental design.} The formal vocabulary the hypothesis reaches for informally is
that of sequential Bayesian experimental design~\cite{rainforth2023boed}, where a design policy maps
history to the next design so as to maximise expected information gain, and can be amortised into a
single forward pass~\cite{foster2021dad}. Our results say something concrete to that framework in
this domain: the expected information gain of a within-run null is close to zero under the surrogate
that is actually available, while the expected information gain of a cross-run null about which
\family to use is large.

\para{Positioning.} Our contribution is the controlled comparison that none of the above performs:
protocol as independent variable, compute held constant, within-family and cross-family adaptivity
separated by an ablated comparator, and the null branch scored for coverage rather than assumed
valid.
